\chapter{Lists and Dictionaries}

\section{Lists}

\subsection{Creation}

\begin{lstlisting}
let nums   = [1, 2, 3, 4, 5]
let states = ["TX", "CA", "NY"]
let mixed  = [1, "two", true]
let empty  = []
\end{lstlisting}

\subsection{Index access}

Indices start at 0. Negative indices count from the end:

\begin{lstlisting}
let xs = [10, 20, 30, 40]
display xs[0]     // 10
display xs[3]     // 40
display xs[-1]    // 40  (last)
display xs[-2]    // 30
\end{lstlisting}

Out-of-bounds access produces a runtime error:

\begin{lstlisting}
xs[10]
// error: index 10 out of bounds (len=4)
\end{lstlisting}

\subsection{Length}

\begin{lstlisting}
len([1, 2, 3])    // 3
\end{lstlisting}

\subsection{Mutation}

\lstinline{push} appends to the end of a list (mutates in place, returns
\lstinline{nil}):

\begin{lstlisting}
let list = [1, 2, 3]
push(list, 4)
display list    // [1, 2, 3, 4]
\end{lstlisting}

\begin{aviso}
  \lstinline{push} is a mutating function — it modifies \lstinline{list}
  directly and returns \lstinline{nil}. Do not write
  \lstinline{let new = push(list, 4)}.
\end{aviso}

\lstinline{pop} removes and returns the last element:

\begin{lstlisting}
let last = pop(list)
\end{lstlisting}

\subsection{Slicing}

\begin{lstlisting}
let xs = [0, 1, 2, 3, 4, 5]
display slice(xs, 1, 4)    // [1, 2, 3]  (indices 1..3)
\end{lstlisting}

\subsection{Iteration}

\begin{lstlisting}
for x in [10, 20, 30] {
  display x * 2
}
\end{lstlisting}

\subsection{Functional operations}

\begin{lstlisting}
let xs = [1, 2, 3, 4, 5]

let doubled  = map(xs, fn(x) { x * 2 })
let evens    = filter(xs, fn(x) { x % 2 == 0 })
let total    = reduce(xs, 0, fn(acc, x) { acc + x })
\end{lstlisting}

\section{Ranges}
\label{sec:ranges}

Ranges are expressions that produce lists of integers. The syntax follows
Rust:

\begin{lstlisting}[caption={Ranges as lists}]
let a = 1..5      // [1, 2, 3, 4]    (exclusive)
let b = 1..=5     // [1, 2, 3, 4, 5] (inclusive)
\end{lstlisting}

Since they are ordinary lists, ranges work with any function that accepts
a list:

\begin{lstlisting}
map(1..=5, fn(x) { x^2 })             // [1, 4, 9, 16, 25]
filter(1..10, fn(x) { x % 2 == 0 })   // [2, 4, 6, 8]
len(1..100)                            // 99
\end{lstlisting}

\subsection{\texttt{chain}}

\lstinline{chain} concatenates multiple lists (or ranges) into one:

\begin{lstlisting}[caption={chain with ranges and lists}]
chain(1..3, 9..=11)           // [1, 2, 9, 10, 11]
chain(1..=3, [10, 20], 5..7)  // [1, 2, 3, 10, 20, 5, 6]
\end{lstlisting}

Typical use: iterating over non-contiguous sequences:

\begin{lstlisting}
for lag in chain(1..=3, 6..=8) {
  display f"lag {lag}"
}
\end{lstlisting}

\subsection{Composition with closures}

Ranges, \lstinline{chain}, and closures compose naturally in a single
expression. The syntax \lstinline{|x| expr} defines a closure without the
ceremony of \lstinline{fn}:

\begin{lstlisting}[caption={Composed expression — single line}]
let r = map(chain(1..=10, 5..9, 56..70), |n| n * 3)
\end{lstlisting}

The pipe also works in \lstinline{for} headers, composing transformations
before entering the loop body:

\begin{lstlisting}[caption={Pipeline in the for header}]
for i in chain(1..=10, 20..=25) |> filter(|n| (n % 2) == 0) {
    display i
}
\end{lstlisting}

\subsection{Comparison with other languages}

\begin{center}
\begin{tabular}{ll}
\toprule
\textbf{Language} & \textbf{Form} \\
\midrule
Python  & \texttt{filter(lambda n: n\%2==0, chain(range(1,11), range(5,9), range(56,70)))} \\
Rust    & \texttt{(1..=10).chain(5..9).chain(56..70).filter(|n| n\%2==0)} \\
Haskell & \texttt{filter even \$ [1..10] ++ [5..8] ++ [56..69]} \\
Julia   & \texttt{Iterators.filter(n->n\%2==0, Iterators.flatten([1:10,5:8,56:69]))} \\
R       & \texttt{Filter(function(n) n\%\%2==0, c(1:10, 5:8, 56:69))} \\
\midrule
Hayashi & \texttt{chain(1..=10,5..9,56..70) |> filter(|n| (n\%2)==0)} \\
\bottomrule
\end{tabular}
\end{center}

\hayashi{} is the only language in which the pipeline runs strictly
left-to-right, with no imports, no namespaces, and a variadic
\lstinline{chain} that treats all arguments symmetrically.

\subsection{Performance note}

\lstinline{chain} materialises ranges into a list before iterating.
\lstinline{for i in 1..N}, by contrast, iterates directly without
allocating. \lstinline{chain} is for small non-contiguous sequences —
lags, years, variable groups. Any contiguous sequence is simply
\lstinline{1..N} or \lstinline{1..=N}.

\section{Dictionaries}

\subsection{Creation}

\begin{lstlisting}
let config = {
  "host": "localhost",
  "port": 5432,
  "ssl": true
}

let empty = {}
\end{lstlisting}

Keys can be strings or integers. Values can be of any type.

\subsection{Access}

\begin{lstlisting}
config["host"]    // "localhost"
config["port"]    // 5432
\end{lstlisting}

A missing key produces an error:

\begin{lstlisting}
config["user"]
// error: key 'user' not found in dict
\end{lstlisting}

Use \lstinline{try}/\lstinline{catch} for safe access:

\begin{lstlisting}
let user = try { config["user"] } catch _ { "anonymous" }
\end{lstlisting}

\subsection{Insertion and update}

\begin{lstlisting}
let d  = {"a": 1}
let d2 = dict_set(d, "b", 2)    // {"a":1, "b":2}
\end{lstlisting}

\lstinline{dict_set} is functional — it returns a new dictionary without
mutating the original.

\subsection{Checking for a key}

\begin{lstlisting}
"host" in config    // true
"password" in config   // false
\end{lstlisting}

\subsection{Dictionaries as function return values}

Analysis functions often return dictionaries with multiple results:

\begin{lstlisting}[caption={summarize returns a dictionary}]
load "data.csv" as df
let stats = summarize(df, income)

display stats["mean"]    // mean
display stats["sd"]      // standard deviation
display stats["N"]       // number of observations
\end{lstlisting}

\subsection{Iterating over dictionaries}

\begin{lstlisting}
let d = {"a": 1, "b": 2, "c": 3}

for key in keys(d) {
  display f"{key} = {d[key]}"
}
\end{lstlisting}
